\begin{derivation}[从\eqnrefs{eq:one_loop_effective_action_general}到\eqnrefs{eq:proper-time-functional-dp11}]

正文\eqnrefs{eq:one_loop_effective_action_general} 表明一圈修正由二次涨落算符的 $\Tr\ln$ 决定。为把紫外问题改写成可按短距离展开的积分，先 Wick 转动到正谱 Euclidean 算符 $\Delta$，再把标量对数恒等式逐个作用在 $\Delta$ 的本征值上。

对正数 $\lambda>0$ 定义
\begin{equation*}
 F(\lambda)
 =-\int_0^\infty\frac{\dd s}{s}
 \left(\ee^{-s\lambda}-\ee^{-s\mu_\kappa^2}\right).
\end{equation*}
对 $\lambda$ 求导，
\begin{equation*}
 \frac{\dd F}{\dd\lambda}
 =\int_0^\infty\dd s\,\ee^{-s\lambda}
 =\frac1\lambda,
\end{equation*}
并且 $F(\mu_\kappa^2)=0$，所以
\begin{equation*}
 F(\lambda)=\ln\frac{\lambda}{\mu_\kappa^2}.
\end{equation*}
若 $\Delta|n\rangle=\lambda_n|n\rangle$，谱定理给出
\begin{align*}
 \Tr\ln\frac{\Delta}{\mu_\kappa^2}
 &=\sum_n\ln\frac{\lambda_n}{\mu_\kappa^2}\\
 &=-\sum_n\int_0^\infty\frac{\dd s}{s}
 \left(\ee^{-s\lambda_n}-\ee^{-s\mu_\kappa^2}\right)\\
 &=-\int_0^\infty\frac{\dd s}{s}
 \Tr\!\left(\ee^{-s\Delta}-\ee^{-s\mu_\kappa^2}\right),
\end{align*}
这就是正文\eqnrefs{eq:proper-time-functional-dp11}。
\end{derivation}

\begin{derivation}[从\eqnrefs{eq:laplace-type-operator-dp11}到\eqnrefs{eq:heat-kernel-a4-flat-dp11}]
在基点 $x$ 选择 Fock--Schwinger 规范

\begin{equation*}
 (y-x)^\mu\omega_\mu(y)=0,
\end{equation*}
则联络在 $x$ 附近为
\begin{equation*}
 \omega_\mu(y)
 =\frac12(y-x)^\nu\Omega_{\nu\mu}(x)
 +O((y-x)^2).
\end{equation*}
自由热核是
\begin{equation*}
 K_0(x,y;s)=\frac1{(4\pi s)^{d/2}}
 \exp\!\left[-\frac{(x-y)^2}{4s}\right].
\end{equation*}
令 $\xi=y-x$，其 Gaussian 矩为
\begin{align*}
 \langle\xi^\mu\xi^\nu\rangle_s
 &=2s\,\delta^{\mu\nu},\\
 \langle\xi^\mu\xi^\nu\xi^\rho\xi^\sigma\rangle_s
 &=4s^2(\delta^{\mu\nu}\delta^{\rho\sigma}
 +\delta^{\mu\rho}\delta^{\nu\sigma}
 +\delta^{\mu\sigma}\delta^{\nu\rho}).
\end{align*}
把 $E$ 与联络项作为固有时有序展开中的插入：一次 $E$ 插入给 $sE$，两次 $E$ 插入给 $s^2E^2/2$；两个线性联络插入结合四阶 Gaussian 矩给
\begin{equation*}
 \frac{s^2}{12}\Omega_{\mu\nu}\Omega^{\mu\nu},
\end{equation*}
而 $E(y)$ 的二阶协变 Taylor 项结合二阶 Gaussian 矩给
\begin{equation*}
 \frac{s^2}{6}\nabla^2E.
\end{equation*}
将这些项与自由热核前因子合并，就得到正文\eqnrefs{eq:heat-kernel-a4-flat-dp11}。
\end{derivation}

\begin{derivation}[从\eqnrefs{eq:heat-kernel-a4-flat-dp11}到\eqnrefs{eq:heat-kernel-divergence-master-dp11}]

对正定欧几里得二次算符 $\Delta$，一圈实玻色 Gaussian 行列式可用正文已经建立的固有时表示写成
\[
\frac12\Tr\ln\Delta
=-\frac12\mu_{\rm DR}^{2\epsilon}
\int_0^\infty\frac{\dd s}{s}\,
\Tr \ee^{-s\Delta},
\]
其中与场无关的加法常数不影响局域反项。把正文\eqnrefs{eq:heat-kernel-a4-flat-dp11} 代入，并只保留含 $s^2$ 的四维局域系数，得到
\begin{align*}
\left.\frac12\Tr\ln\Delta\right|_{a_4}
={}&-\frac12\mu_{\rm DR}^{2\epsilon}
\int\dd^d x\,
\tr a_4(x)\\
&\times\int_0^\infty\frac{\dd s}{s}
\frac{s^2}{(4\pi s)^{d/2}}.
\end{align*}
取 $d=4-2\epsilon$ 后，短固有时端的幂次变为
\[
\frac1s\frac{s^2}{s^{2-\epsilon}}
=s^{-1+\epsilon}.
\]
因此紫外极点只需从 $s\to0^+$ 端读取。任取固定分界尺度 $s_0$，
\begin{align*}
\mu_{\rm DR}^{2\epsilon}
\int_0^{s_0}\dd s\,s^{-1+\epsilon}
&=\mu_{\rm DR}^{2\epsilon}
\frac{s_0^\epsilon}{\epsilon}\\
&=\frac1\epsilon
+\ln(\mu_{\rm DR}^2s_0)+O(\epsilon).
\end{align*}
所以极点系数与任意分界 $s_0$ 无关，而长固有时部分只改变有限项。代入
\[
 a_4(x)
 =\frac12E^2
 +\frac1{12}\Omega_{\mu\nu}\Omega^{\mu\nu}
 +\frac16\nabla^2E
\]
并使用无边界体积分
\(
\int\dd^4 x\,\nabla^2E=0
\)，得到
\[
\left.\frac12\Tr\ln\Delta\right|_{\rm pole}
=-\frac1{2(4\pi)^2}\frac1\epsilon
\int\dd^4 x\,
\tr\left[
\frac12E^2+\frac1{12}\Omega_{\mu\nu}\Omega^{\mu\nu}
\right].
\]
最后使用正文已经固定的 $1/\bar\epsilon$ 记号整理 $1/\epsilon-\gamma_E+\ln4\pi$，即得到\eqnrefs{eq:heat-kernel-divergence-master-dp11}。
\end{derivation}
